%% requires: \providecommand{\vfy}{\textsuperscript{[v]}}  (marks a value still to be verified)
% ============================================================
% GROUP 5: Climate hazard exposure (computed in F4, hazard_context)
% Place: Literature Review (section 1) and Methodology (section 2).
% Scope decision: exposure indicators only. No loss functions, no derating, no cut-out logic.
% ============================================================

\section{Literature Review: Climate Hazards for Photovoltaic and Wind Plants}
\label{sec:g5-lit}

\subsection{Operating envelopes in standards}
Wind turbine design standards define the ambient conditions a machine must withstand. IEC 61400-1 sets a normal ambient temperature range of $-10$ to $+40\,^{\circ}\mathrm{C}$ for operation and an extreme range of $-20$ to $+50\,^{\circ}\mathrm{C}$, and defines extreme wind and gust models \citep{iec61400}. The 2025 interpretation sheet of the 2019 edition clarifies that the turbine must be able to operate and generate up to $+40\,^{\circ}\mathrm{C}$ without being required to deliver rated power at that temperature, so output derating above $40\,^{\circ}\mathrm{C}$ is permitted by the standard. % verify: subclauses 6.4.2 and 6.4.3.2 of the 2019 edition were read in the interpretation sheet and in Indian grid authority minutes, not in the full text of the edition
India's grid authority has moved in the opposite direction, requiring new turbines not to derate below $40\,^{\circ}\mathrm{C}$ for grid connectivity. % verify: reading of the minutes of 26 July 2024
For this reason a count of days above $40\,^{\circ}\mathrm{C}$ is a physically meaningful exposure indicator for wind, because it marks the upper edge of the normal operating range, without implying a specific loss.

\subsection{Why losses are not modeled}
Turning exposure into energy loss requires model-specific curves: the temperature coefficient of the PV module, the derating onset and slope of the inverter, and the shutdown temperature. For inverters these curves are manufacturer and model specific and, in the European standard for inverter efficiency, the derating curve is optional. % verify: EN 50524:2021 statement taken from a conference paper on inverter derating
No set of primary datasheets for utility-scale inverters was assembled for this work, and any loss function built without them would add Tier~3 parameters to the economic result. The framework therefore reports hazard exposure next to the robust ranking and does not convert it into LCOE.

\subsection{Insulation and moisture}
The safety standard for photovoltaic inverters does not fix a single insulation resistance: the minimum DC insulation resistance of the array scales with the maximum PV voltage, $R = V_{\max} / 30\,\mathrm{mA}$ \citep{iec62109}, which is about $33\,\mathrm{k\Omega}$ at 1000~V and $50\,\mathrm{k\Omega}$ at 1500~V. % verify: formula read from a test report form (clauses 4.8.2.1 and 4.8.2.2), not from the standard itself
A fixed $500\,\mathrm{k\Omega}$ threshold found in a manufacturer article has no normative basis and is not used. A binary flood or moisture loss based on it is not part of the framework.

\subsection{Projected change in wind resource}
Studies of wind under climate change report large differences between models and regions. A study of 14 downscaled CMIP6 models covers low-production extremes globally but reports no regional ratios of mean wind speed. % verify: paper not fully referenced
For South America, a regional study reports changes in wind power density from more than $+100\%$ in parts of Brazil, Paraguay, Uruguay and Venezuela to about $-50\%$ along the west coast between 1995--2014 and 2081--2100. \citep{martinez2024} % verify: only the abstract page was seen
For Spain, four wind farms were studied with EURO-CORDEX regional simulations under CMIP5 scenarios until 2065 \citep{solaun2020}, and a review compiles quantitative projections of renewable generation under climate change \citep{solaun2019}. No source found for this work gives numerical ratios of future to historical mean wind speed by grid cell or percentile for South America, the Iberian Peninsula and South Asia under CMIP6. The plausible range of the wind factor in the framework is therefore not taken from the literature.

\section{Methodology: Hazard Exposure Indicators}
\label{sec:g5-method}

\subsection{Indicators}
Exposure is computed per grid cell and per climate scenario, for the reference period and for future windows, and is reported next to the regret ranking. It does not enter the LCOE and it does not enter the regret.

\begin{table}[ht]
\centering
\caption{Hazard exposure indicators.}
\label{tab:g5-ind}
\small
\begin{tabular}{p{2.6cm}p{5.2cm}p{1.9cm}p{2.9cm}}
\toprule
Indicator & Definition & Tier & Use \\
\midrule
\texttt{tx35\_days} & Days per year with daily maximum 2~m temperature above $35\,^{\circ}\mathrm{C}$ & 3 (author threshold) & exposure, PV and wind \\
\texttt{tx40\_days} & Days per year above $40\,^{\circ}\mathrm{C}$ & 1 for wind (edge of the normal operating range), 3 for PV & exposure \\
\texttt{rx5day\_mm} & Annual maximum of 5-day accumulated precipitation, averaged over the years of the window & 3 & exposure \\
\texttt{wet\_p95\_exceed\_freq} & Share of wet days (precipitation of at least 1.0~mm) above the cell's own 95th percentile of wet-day precipitation in the reference period 1995--2014 & 3 & exposure \\
\texttt{gust\_mean\_annual\_max\_ms} & Mean of the annual maximum 10~m wind gust from ERA5, 0.25$^{\circ}$, 1995--2014 & n/a & context only, not scored \\
\bottomrule
\end{tabular}
\end{table}

The gust indicator is invariant to the climate scenario and has no threshold. The ERA5 gust is a near-surface (10~m) value, whereas turbine design wind speeds refer to hub-height averages and to a 50-year return period. A comparison with a turbine cut-out or design limit would require extrapolation in height and in return period, which is not performed. The threshold remains undefined until the open question on the gust channel is closed. % open: OQ-051

\subsection{Validity range of the wind factor}
Scenario wind multipliers outside a validity range are rejected as implausible. The current range is 0.5 to 1.5. This range is a Tier~3 guard and is not derived from published ratios. The share of Brazilian cell-members it masks was examined as a sensitivity: 1.01\% at a range of 0.4 to 1.6, 1.43\% at the current range of 0.5 to 1.5 (members MIROC6, GFDL-ESM4 and IPSL-CM6A-LR) and 4.3\% at 0.7 to 1.3. In Portugal no cell-member is masked down to 0.8 to 1.2, and in India at most 0.2\% are. % D-F4-018 One ensemble member (IPSL-CM6A-LR) reaches a mean ratio of 1.198 in 2071--2100 under SSP5-8.5 in the project's own output, inside the range on average, although 17,802 Brazilian cell-members of that member are masked by it. % verify: confirm run identifier and value
The member is not excluded: it is run with and without, and the effect on the ranking is reported. As a rough check, the cube root of a $+25\%$ to $+50\%$ change in wind power density corresponds to a speed ratio of about 1.08 to 1.15. This is the author's conversion and not a published result. % verify: author calculation, not a literature value

\subsection{Limitations specific to Group 5}
\begin{itemize}
\item Exposure is not converted into loss; hazard and economic results are read together.
\item The 35 and $40\,^{\circ}\mathrm{C}$ thresholds are simple counts; only the latter has a physical anchor for wind.
\item The wind-factor range is a Tier~3 guard; no CMIP6 per-cell ratio was found for the study regions.
\item Gust exposure is context only and is not comparable with design wind speeds.
\end{itemize}
